\input{\string~/Documents/School/"UC Irvine"/ta_template2.0.tex}
\rh{2B}{5.1-6.5}
\chead{\today}
\graphicspath{ {\string~/Desktop} }
\usetikzlibrary{arrows}
%============ANSWER KEY TOGGLE===========
\newcommand{\ans}[1]{ \noindent {\color{red} \textbf{Answer:} #1}}
%\renewcommand{\ans}[1]{\vspace{3in}}
%==========================================
\begin{document}
\begin{center}
\textbf{Midterm Practice Problems}
\end{center}

\ans{The source of these problems, and more like them, is \href{https://www.shorturl.at/ktyNY}{here}}

\problembox{[Challenge] Evaluate the limit:
\[
\lim_{n \to \infty} \sum_{i = 1}^n \lrp{\frac{i^4}{n^5} + \frac{i}{n^2}}
\]
}

\ans{The trick here is to evaluate this as an integral. First, factor out $1/n$ from the summand:
\[
\frac{i^4}{n^5} + \frac{i}{n^2} = \frac1n \lrp{ \lrp{i/n}^4 + (i/n)}
\]
from here, we see that $a = 0$, $\Delta x = 1/n$ so $b = 1$. We're home free when we write down:
\[
\int_0^1 x^4 + x dx = \lrb{\frac{x^5}{5} + \frac{x^2}{2}}_0^1 = \frac{1}{5} + \frac12 = \frac{7}{10}
\]
}

\problembox{Compute the following antiderivative:
\[
\int \cos^2 (x) - \sin^2(x) dx
\]
}
\ans{Use difference of squares and then $u$ substitution}

\problembox{The area of the region bounded by the curve $y=e^{2 x}$, the $x$-axis, the $y$-axis, and the line $x=2$ }

\ans{C: $\frac{e^4}{2}-\frac{1}{2}$ square units
\[
\begin{aligned}
A & =\int_0^2 e^{2 x} d x=\frac{1}{2} \int_0^2 e^{2 x}(2) d x \\
& =\frac{1}{2}\left[e^{2 x}\right]_0^2 \\
& =\frac{1}{2}\left(e^4-e^0\right) \\
& =\frac{e^4-1}{2} \\
& =\frac{e^4}{2}-\frac{1}{2}
\end{aligned}
\]
}


\problembox{The average value of $y=e^{3 x}$ over the interval from $x=0$ to $x=4$ is}

\ans{
\[
\text { average value } \begin{aligned}
& =\frac{1}{b-a} \int_a^b f(x) d x \\
& =\frac{1}{4-0} \int_0^4 e^{3 x} d x \\
& =\frac{1}{4} \cdot \frac{1}{3} \int_0^4 e^{3 x} \cdot 3 d x \\
& =\frac{1}{12}\left[e^{3 x}\right]_0^4 \\
& =\frac{1}{12}\left(e^{12}-e^0\right) \\
& =\frac{e^{12}-1}{12}
\end{aligned}
\]
}

\problembox{Which of the following is equivalent to
$$
\lim _{n \rightarrow \infty} \sum_{i=1}^n\left[\left(1+\frac{2 i}{n}\right)^2+1\right]\left(\frac{2}{n}\right) ?
$$\\
A. $\int_2^4\left(x^2+1\right) d x$ \\
B. $\int_1^3 x^2 d x$ \\
C. $\int_1^2\left(x^2+1\right) d x$ \\
D. $\int_1^3\left(x^2-1\right) d x$ \\
E. $\int_1^3\left(x^2+1\right) d x$ \\}

\ans{E}

\problembox{13. Let $A$ be the area bounded by one arch of the sine curve. Which of the following represents the volume of the solid generated when $A$ is revolved around the $x$-axis? \\
A. $2 \pi \int_0^\pi x \sin x d x$\\
B. $\pi \int_0^\pi \sin ^2 x d x$\\
C. $\pi \int_0^\pi x \sin x d x$\\
D. $\pi \int_0^{2 \pi} \sin ^2 x d x$\\
E. $2 \pi \int_0^1 \arcsin y d y$}

\ans{B}


\problembox{16. Approximate the value of $\int_1^3 \ln x d x$ using 4 circumscribed rectangles.\\
A. $1.007$\\
B. $1.296$\\
C. $1.557$\\
D. $2.015$\\
E. $3.114$}
\ans{C
\[
\begin{aligned}
\text { Area } & \approx \frac{1}{2} \ln \frac{3}{2}+\frac{1}{2} \ln 2+\frac{1}{2} \ln \frac{5}{2}+\frac{1}{2} \ln 3 \\
& =\frac{1}{2}\left(\ln \frac{3}{2}+\ln 2+\ln \frac{5}{2}+\ln 3\right) \\
& \approx 1.557
\end{aligned}
\]
}
\end{document}